% Generated draft: requires mathematical and source review.
\documentclass[11pt]{article}
\usepackage[margin=25mm]{geometry}
\usepackage{fontspec}
\setmainfont{Latin Modern Roman}
\usepackage{amsmath,amssymb,mathtools,unicode-math}
\setmathfont{Latin Modern Math}
\usepackage{enumitem,needspace,xurl,hyperref}
\setlength{\parindent}{0pt}
\setlength{\parskip}{5pt}
\setlength{\emergencystretch}{3em}
\newcommand{\N}{\mathbb{N}}
\newcommand{\Z}{\mathbb{Z}}
\newcommand{\Q}{\mathbb{Q}}
\newcommand{\R}{\mathbb{R}}
\newcommand{\C}{\mathbb{C}}
\newcommand{\F}{\mathbb{F}}
\newcommand{\A}{\mathbb{A}}
\newcommand{\PP}{\mathbb{P}}
\newcommand{\E}{\mathbb{E}}
\newcommand{\eps}{\varepsilon}
\newcommand{\dd}{\,\mathrm{d}}
\begin{document}
\section*{UnsolvedMath ID: 700255}
\textbf{Source problem code:} catalogue entry
\subsection{Definitions and mathematical statement}
For a prime \(p\), a \(p\)-subgroup of a group is a subgroup whose elements all have finite order equal to a power of \(p\). A Sylow \(p\)-subgroup means a subgroup maximal among the \(p\)-subgroups. A group \(G\) is locally finite if every finitely generated subgroup of \(G\) is finite. Following the supplied source, an \(A\)-group is a locally finite group all of whose Sylow \(p\)-subgroups, for every prime \(p\), are abelian.

A group identity \(w(x_1,\ldots,x_n)=1\) is an equation given by a group word that holds under every substitution of group elements for its variables. A variety of groups is the class of all groups satisfying some set of group identities. Here a variety of \(A\)-groups is understood to mean a variety every member of which is an \(A\)-group; in particular, every member is required to be locally finite.

\textbf{Problem (A. Yu. Olshanskii, 1973; Kourovka Notebook 4.48).} Does every variety \(\mathcal V\) of \(A\)-groups have a finite basis of identities? Explicitly, for every such \(\mathcal V\), do there exist finitely many group words \(w_1,\ldots,w_m\), each in finitely many variables, such that for every group \(H\),
\[
H\in\mathcal V
\quad\Longleftrightarrow\quad
H\models w_i=1\ \text{for every }1\leq i\leq m?
\]
The equivalence is required among all groups, so local finiteness and the Sylow conditions are hypotheses on the members of \(\mathcal V\), rather than extra restrictions on the test group \(H\).
\subsection{Short English statement}
If every group in a group variety is locally finite and has abelian Sylow subgroups, must that variety be definable by finitely many group identities? The supplied registry labels this problem open; that status has not been independently verified.
\subsection{Sources}
\begin{itemize}
\item Supplied local UnsolvedMath registry record, permanent ID \(700255\), exact source problem code \texttt{catalogue entry}, titled \emph{Kourovka Notebook Problem 4.48}. The record supplies the status \texttt{open} and an update date of 2026-10-16; these are registry metadata, not an independently verified literature assessment.
\item \emph{The Kourovka Notebook (Unsolved Problems in Group Theory)}, Problem 4.48, attributed to A. Yu. Olshanskii, proposed in 1973. The supplied background identifies the 4th issue (1973) and the 21st issue edition (2026), arXiv:1401.0300v48.
\item The supplied collection description identifies E. I. Khukhro and V. D. Mazurov as editors of the Kourovka Notebook collection covering Issues 1--21 (1965--2026). No source was visited or independently checked for this draft.
\end{itemize}
\end{document}
